\begin{Exercise}[label={ex:archivio}]
    Un file di testo "archivio.txt" contiene i dati (primo autore, titolo, numero di copie possedute, 
    numero di copie in prestito) relativi ai differenti volumi conservati presso una biblioteca.
    Più precisamente, ogni riga del file contiene nell'ordine, separati da uno spazio bianco:
    \begin{itemize}
        \item autore (al più 23 caratteri senza spazi)
        \item titolo (al più 55 caratteri senza spazi)
        \item numero di copie possedute (intero)
        \item numero di copie in prestito (intero)
    \end{itemize}

    Si realizzi un programma C che:
    \begin{enumerate}
        \item Legga il contenuto di "archivio.txt" e costruisca in memoria un vettore V di strutture
        corrispondenti (si supponga che il file "archivio.txt" non possa contenere più di 30 righe). 
        Si stampi a video il contenuto del vettore.
        \item A partire da V, costruisca una lista L di interi contenente per ciascun volume il 
        numero di copie disponibili nella biblioteca, ovvero la differenza fra il numero di copie 
        possedute e il numero di copie in prestito. Si stampi a video il contenuto della lista L.
        \item Utilizzando L per ottenere la somma delle copie disponibili e V per la somma delle
        copie possedute, calcoli il rapporto fra volumi disponibili e volumi posseduti.
        \item  Utilizzando la lista di interi L, stampi il numero di riga di "archivio.txt" 
        relativo al volume con più copie disponibili. In caso di più volumi con pari numero di copie 
        disponibili, qualunque riga relativa a questi ultimi è considerata una risposta corretta.
    \end{enumerate}
\end{Exercise}

\begin{Answer}[ref={ex:archivio}]
  	\lstinputlisting[style=cstylesmall]{src/archivio.c}
\end{Answer}